% TOP_PROBLEM: {"id": 20000925, "title": "The exceptional three-color path problem: star-forest and triangle-density reductions", "category_id": 2, "category": {"id": 2, "name": "combinatorics", "display_name": "Combinatorics", "description": "Counting problems, graph theory, discrete structures.", "slug": "combinatorics", "order_index": 2, "created_at": "2026-07-31T15:26:25.671Z"}, "status": "partially_solved", "problem_number": "AIM-COMBINATORICS-0050", "set_id": 14}
% FIRST_POSTED: 2026-10-01
% ENTRY_KIND: statement_only
% UnsolvedMath ID 20000925; source catalogue code AIM-COMBINATORICS-0050
% !TeX program = lualatex
% Definition and sources only; this file is not an LLM solution attempt.
\documentclass[11pt,a4paper]{article}
\usepackage[margin=25mm]{geometry}
\usepackage{fontspec}
\setmainfont{lmroman10-regular.otf}[BoldFont=lmroman10-bold.otf,
 ItalicFont=lmroman10-italic.otf,BoldItalicFont=lmroman10-bolditalic.otf]
\usepackage{amsmath,amssymb,mathtools}
\usepackage{unicode-math}
\setmathfont{latinmodern-math.otf}
\AtBeginDocument{\let\setminus\smallsetminus}
\usepackage{xurl,hyperref}
\hypersetup{colorlinks=true,urlcolor=blue}
\setcounter{secnumdepth}{0}
\setlength{\parindent}{0pt}
\setlength{\parskip}{5pt}
\setlength{\emergencystretch}{3em}
\begin{document}
\section*{The exceptional three-color path problem: star-forest and triangle-density reductions}\label{20000925}
{\small UnsolvedMath ID 20000925 · AIM-COMBINATORICS-0050 · Combinatorics · AIM Workshop Problem Lists}
\subsection{Definitions and mathematical statement}
A graph here is finite, simple, and undirected. A proper vertex coloring
assigns different colors to adjacent vertices; the chromatic number
$\chi(G)$ is the minimum number of colors permitting such an assignment.
An edge coloring with three colors is an arbitrary function
$c:E(G)\to\{1,2,3\}$ and is not required to be proper.
A path of length three consists of four distinct vertices
$v_0,v_1,v_2,v_3$ with edges $v_0v_1,v_1v_2,v_2v_3$.
It is monochromatic when these three edges receive the same color.
Extra edges between its vertices are permitted.

Gy\'arf\'as's question is whether
\[
 \chi(G)\ge6\quad\Longrightarrow\quad
 \text{every }c:E(G)\to\{1,2,3\}
 \text{ has a monochromatic path of length three}.
\]
The universal quantifiers range over all such graphs and all three-color
assignments to their edges. Equivalently, suppose the edge set can be
partitioned into three spanning subgraphs, each containing no
three-edge path. Must their union be properly vertex colorable using
at most five colors? A missing edge belongs to none of the three color
classes and cannot be used to form a path. Thus the claim concerns
arbitrary graphs of chromatic number at least six, without any
completeness or minimum-degree assumption. The path notation in the
workshop counts edges, so its $P_3$ is the four-vertex path.
\subsection{Short English statement}
Does every three-coloring of the edges of a graph requiring six vertex colors contain a monochromatic three-edge path?
\subsection{Sources}
\begin{itemize}
\item[S1] ULAM.AI and the UnsolvedMath Contributors, supplied problems.json, record 20000925 (AIM-COMBINATORICS-0050), AIM Workshop Problem Lists. Catalogue identity and source transcription; CC BY 4.0.
\url{https://huggingface.co/datasets/ulamai/UnsolvedMath}.
\item[S2] American Institute of Mathematics, Graph Ramsey theory, item 3. Original workshop question underlying AIM-COMBINATORICS-0050.
\url{https://www.overleaf.com/read/mnvcscjjysvg}.
\end{itemize}
\end{document}
